\documentclass[../../../include/open-logic-section]{subfiles}

\begin{document}

\olfileid[id]{his}{set}{cantorplane}
\olsection{Cantor tentang Garis dan Bidang}

Beberapa keadaan seputar pembuktian Schr\"oder--Bernstein berkaitan dengan
sejarah yang kita bahas dalam \olref[his][set][pathology]{sec}. Ingat bahwa
pada 1877 Cantor membuktikan bahwa terdapat tepat sama banyak titik pada
suatu persegi dengan pada salah satu sisinya. Di sini kita akan menyajikan
bukti (percobaan pertama) Cantor.

Misalkan $\unitline$ garis satuan, yakni himpunan titik $[0,1]$. Misalkan
$\unitsquare$ persegi satuan, yakni himpunan titik
$\unitline\times\unitline$. Dalam istilah ini, Cantor membuktikan
$\cardeq{\unitline}{\unitsquare}$. Ia menulis suatu catatan kepada Dedekind
yang pada dasarnya memuat argumen berikut.

\begin{thm}\ollabel{thm:cantorplane}
$\cardeq{\unitline}{\unitsquare}$.
\end{thm}

\begin{proof}[Bukti: bagian pertama.]
Tetapkan $a,b\in\unitline$. Tuliskan keduanya dalam notasi biner, sehingga
kita mempunyai barisan tak hingga angka $0$ dan $1$, yaitu
$a_1,a_2,\dots$ dan $b_1,b_2,\dots$, sedemikian sehingga:
\begin{align*}
a &= 0.a_1a_2a_3a_4\dots\\
b &= 0.b_1b_2b_3b_4\dots
\intertext{Sekarang pertimbangkan fungsi
$f \colon \unitsquare \to \unitline$ yang diberikan oleh}
f(a, b) & = 0.a_1b_1a_2b_2a_3b_3a_4b_4\dots
\end{align*}
Sekarang $f$ merupakan !!a{injection}, sebab jika $f(a,b)=f(c,d)$, maka
$a_n=c_n$ dan $b_n=d_n$ bagi semua $n\in\Nat$, sehingga $a=c$ dan $b=d$.
\end{proof}

Sayangnya, sebagaimana ditunjukkan Dedekind kepada Cantor, hal ini tidak
menjawab pertanyaan semula. Pertimbangkan
$0.\dot{1}\dot{0}=0.1010101010\ldots$. Kita memerlukan
$f(a,b)=0.\dot{1}\dot{0}$, dengan:
\begin{align*}
a&= 0.\dot{1}\dot{1} = 0.111111\ldots\\
b&= 0
\end{align*}
Namun, $a=0.\dot{1}\dot{1}=1$. Jadi, ketika kita berkata
``tuliskan $a$ dan $b$ dalam notasi biner'', kita harus memilih
\emph{notasi mana} yang digunakan; dan karena $f$ harus merupakan suatu
\emph{fungsi}, kita hanya dapat menggunakan \emph{satu} dari kedua notasi
yang mungkin. Namun, jika misalnya kita menggunakan notasi sederhana dan
menulis $a$ sebagai ``$1.000\ldots$'', tidak ada pasangan
$\tuple{a,b}$ sedemikian sehingga
$f(a,b)=0.\dot{1}\dot{0}$.

Ringkasnya, Dedekind menunjukkan bahwa karena kemungkinan adanya ekspansi
biner berulang tertentu, fungsi $f$ dari Cantor merupakan !!a{injection}
tetapi \emph{bukan} !!a{surjection}. Jadi, Cantor hanya menunjukkan
$\cardle{\unitsquare}{\unitline}$ dan \emph{bukan}
$\cardeq{\unitsquare}{\unitline}$.

Cantor segera membalas Dedekind dan pada dasarnya mengemukakan bahwa bukti
dapat diselesaikan sebagai berikut:

\begin{proof}[Bukti: penyelesaian.]
Kita telah menunjukkan $\cardle{\unitsquare}{\unitline}$. Namun, jelas
terdapat !!a{injection} dari $\unitline$ ke $\unitsquare$: cukup letakkan
garis secara mendatar pada salah satu sisi persegi. Jadi,
$\cardle{\unitline}{\unitsquare}$ dan
$\cardle{\unitsquare}{\unitline}$. Berdasarkan
Schr\"{o}der--Bernstein
(\olref[sfr][siz][sb]{thm:schroder-bernstein}),
$\cardeq{\unitline}{\unitsquare}$.
\end{proof}

Namun, tentu saja Cantor tidak dapat menyelesaikan baris terakhir dengan
istilah-istilah ini karena Teorema Schr\"{o}der--Bernstein belum terbukti.
Memang, walaupun kemudian Cantor merumuskannya sebagai dugaan umum, teorema
itu baru dibuktikan secara memuaskan pada 1897. (Karena itu, kemudian pada
1877 Cantor memberikan bukti lain bagi \olref{thm:cantorplane}, yang tidak
melalui Schr\"{o}der--Bernstein.)

\end{document}
